\chapter{Peripheral Channel Structure and Transfer-Operator Gaps}\label{ch:spectral}

The spectral theory of finite-dimensional quantum channels and MPS transfer maps
rests on two principal structures, following Pérez-García et
al.~\cite{PerezGarcia2007Matrix} and Wolf~\cite{Wolf2012Quantum}: the peripheral
spectrum and the mixed transfer operator.

First, for an irreducible unital Schwarz map with a faithful adjoint fixed
point, the peripheral eigenvalues on the unit circle form a finite cyclic group
(Theorem~\cite[``Peripheral eigenvalues: cyclic structure for Kraus
    maps'']{QICLean2026Blueprint}). Under trace preservation, the order of this
group divides the bond dimension~$D$ (Theorem~\cite[``Transfer-map peripheral
    eigenvalues form a cyclic group'']{QICLean2026Blueprint}). Unitary
conjugation and cyclic corner projections yield the periodicity-free
cyclic-sector decomposition (Theorem~\cite[``Cyclic decomposition for
    irreducible finite Kraus maps'']{QICLean2026Blueprint}) used in the
canonical-form reduction of Chapter~\ref{ch:canonical}.

Second, the \emph{mixed transfer operator} $F_{AB}$ of two MPS tensors $A$ and
$B$ governs the asymptotic overlap of matrix product vectors. Under
trace-preserving normalization, the spectral radius satisfies
$\rho_{\spec}(F_{AB})\leq1$ (Theorem~\ref{thm:eigenvalue_bound}); for injective
(or irreducible trace-preserving) tensors, equality $\rho_{\spec}(F_{AB})=1$
forces $A$ and $B$ to be gauge-phase equivalent
(Theorem~\ref{thm:modulus_one_gauge}), while unequal bond dimensions force a
strict gap $\rho_{\spec}(F_{AB})<1$
(Theorem~\ref{thm:transfer_operator_gap_rect}). Because the MPV overlap
$O_{AB}(N)$ equals the operator trace $\Tr(F_{AB}^N)$
(Theorem~\ref{thm:overlap_trace}), these spectral gaps imply the overlap-decay
theorems (Theorem~\ref{thm:overlap_decay} and its rectangular and
irreducible-trace-preserving variants). For a single primitive tensor whose
transfer map has trivial peripheral spectrum away from the fixed-point
projection, the complementary gap yields the rank-one Perron limit
$\Tr(\E_A^n)\to1$ (Theorem~\ref{thm:trace_pow_one}) and primitive self-overlap
convergence $O_{AA}(N)\to1$ (Theorem~\ref{thm:primitive_overlap}).

Throughout, $\rho_{\spec}(\cdot)$ denotes the spectral radius of a linear
endomorphism of a finite-dimensional matrix space: $\rho_{\spec}(F_{AB})<1$ for
a mixed transfer operator and $\rho_{\spec}(\E_A-P)<1$ for a transfer map after
subtracting its fixed-point projection $P$. The unqualified phrase \emph{spectral
    gap} is reserved for the energy gap of a many-body Hamiltonian.

The algebraic mixed-transfer identities hold without normalization. The
subsequent spectral-radius estimates assume the trace-preserving normalization
\begin{align}
    \sum_i (A^i)^\dagger A^i
     & =\Id.
    \label{eq:spec_tp_normalization}
\end{align}
Right-canonical (unital) and left-canonical (trace-preserving) gauges are
distinct similarity transformations of a tensor (Remark~\ref{rem:gauge_conventions}),
used according to whether the fixed point of the transfer map or that of its adjoint
governs the transformation.

Supporting proofs for the mixed-transfer identities, operator-norm estimates,
gauge intertwiners, and cyclic-decomposition details appear in
Appendix~\ref{ch:spectral_supporting_results}.
